\documentclass[11pt]{article}
\usepackage[letterpaper,margin=1in,headheight=15pt]{geometry}
\usepackage[T1]{fontenc}
\usepackage{lmodern}
\usepackage{amsmath,amssymb,amsthm,mathtools}
\usepackage{booktabs,longtable,array,enumitem,microtype,xcolor,xurl}
\usepackage{fancyhdr}
\usepackage[colorlinks=true,linkcolor=blue!45!black,citecolor=blue!45!black,urlcolor=blue!45!black]{hyperref}
\usepackage{bookmark}
\setlength{\emergencystretch}{3em}
\setcounter{tocdepth}{2}
\numberwithin{equation}{section}
\newtheorem{theorem}{Theorem}[section]
\newtheorem{proposition}[theorem]{Proposition}
\newtheorem{lemma}[theorem]{Lemma}
\newtheorem{corollary}[theorem]{Corollary}
\theoremstyle{definition}
\newtheorem{definition}[theorem]{Definition}
\newtheorem{example}[theorem]{Example}
\theoremstyle{remark}
\newtheorem{remark}[theorem]{Remark}
\DeclareMathOperator{\Li}{Li}
\DeclareMathOperator{\Log}{Log}
\DeclareMathOperator{\Cl}{Cl}
\DeclareMathOperator{\NF}{NF}
\DeclareMathOperator{\rank}{rank}
\DeclareMathOperator{\Tor}{Tor}
\DeclareMathOperator{\ord}{ord}
\DeclareMathOperator{\im}{im}
\DeclareMathOperator{\coker}{coker}
\DeclareMathOperator{\cond}{cond}
\DeclareMathOperator{\tr}{tr}
\newcommand{\Z}{\mathbb Z}
\newcommand{\Q}{\mathbb Q}
\newcommand{\C}{\mathbb C}
\newcommand{\F}{\mathbb{F}_2}
\newcommand{\U}{\mathcal U}
\newcommand{\I}{\mathcal I}
\newcommand{\B}{\mathcal B}
\newcommand{\Lcal}{\mathcal L}
\newcommand{\HH}{\widehat H}
\newcommand{\eps}{\varepsilon}
\newcommand{\code}[1]{\texttt{\small #1}}
\newcommand{\src}[1]{\nolinkurl{#1}}
\newcommand{\snap}{28357e8ca63dd78327db91d9be239d75e4462879}
\pagestyle{fancy}
\fancyhf{}
\fancyhead[L]{\small Integral distributions and reflection torsion}
\fancyhead[R]{\small ProveIt continuation}
\fancyfoot[C]{\thepage}
\hypersetup{pdftitle={Unimodular Distribution Certificates and Exact Reflection Torsion for Cyclotomic Polylogarithms},pdfauthor={Research report prepared for Vladimir Reshetnikov}}

\title{\textbf{Unimodular Distribution Certificates\\and Exact Reflection Torsion\\for Cyclotomic Polylogarithms}\\[6pt]
\large Integral normal forms, arbitrary characteristic, and finite-order jets}
\author{Research report prepared for Vladimir Reshetnikov\\\small A continuation of the ProveIt Polylogarithms manuscript}
\date{9 October 2026}
\begin{document}
\maketitle
\begin{abstract}
The characteristic-zero distribution normal form in the ProveIt manuscript
leaves integral lattices, modular specialization, and reflection torsion as
separate questions. We give a fixed basis consisting of original grid
symbols and a unit-determinant minor made from unscaled prime distribution
rows. Consequently the complete weighted relation matrix is unimodularly
equivalent to an identity block and zeros over
$\Z[A_p:p\mid q]$, and remains so after arbitrary commutative-ring base
change. A division-free recurrence supplies exact polynomial certificates
and explicit coefficient bounds.

Reflection behaves differently. A weighted Anderson resolution reduces
its Tate cohomology to a finite Koszul complex over the coefficient ring
modulo two. This gives the complete Smith torsion for all integer prime
weights. For weights in $\Z[u]/(u^L)$, the answer is controlled by one
minimum binary valuation: the torsion multiplicity is
$v\,2^{r_*-1}$, where $r_*$ counts the odd prime factors and also counts
$2$ precisely when $4\mid q$. We also exhibit a reflected jet module whose
torsion-free part is not free over the jet ring.

Specializing the normal forms to $A_p=p^{1-s}$ produces meromorphic
polylogarithm identities for every complex order, including explicit
level-$12$ and level-$30$ formulas and pole-cancelled identities for all
order derivatives at $s=1$. The package includes independent exact
certificate replay, direct Smith computations, and separately labelled
numerical diagnostics. The ordinary-distribution and resolution
antecedents of Kubert, Anderson, and Ouyang are retained; no worldwide
priority or numerical period-independence claim is made.
\end{abstract}
\vspace{5pt}
\noindent\textbf{Scope.} All theorem and proposition statements below have
ordinary mathematical proofs. The executable checks test finite instances
and implementations; they are not proof-assistant formalizations. The
remaining mixed-depth $S_6$ identity is not proved in this report.
\clearpage
\tableofcontents
\clearpage

\section{Research target and relation to the existing manuscript}
\label{sec:scope}

\subsection{What is being completed}
The input manuscript \emph{Polylogarithms and their Arithmetic Bridges}
and its distribution-jets report are read at the repository snapshot
\begin{center}
\small\texttt{\snap}.
\end{center}
The report's universal conductor normal form is stated over a
characteristic-zero field containing character values. Its polynomial
weights may be specialized in any algebra over that field. The manuscript
explicitly leaves integral Smith forms, torsion, and integral certificate
arithmetic for further work \cite{repo,repodistribution}.

The distinction is substantive. A Fourier basis change divides by group
orders. Thus a characteristic-zero proof of rank does not, by itself,
prove saturation over $\Z$, preserve a basis in characteristic two, or
identify the torsion introduced by the equations $e_{-x}=\pm e_x$.
Conversely, the vanishing of a determinant attached to a particular
primitive-grid coordinate map need not be a failure of the complete
presentation. This report treats the complete prime-row presentation,
not an unspecified numerical lattice of periods.

The principal results are the original-symbol unit-minor theorem
(Theorem~\ref{thm:unimodular}), the reflection Koszul reduction
(Theorem~\ref{thm:koszul}), and the exact one-variable jet formula
(Theorem~\ref{thm:jets}). Together they settle these integral and modular
questions for the stated module, including its reflected scalar and
one-variable-jet specializations. The polylogarithm consequences are
proved in Section~\ref{sec:polylog}.

\subsection{Classical antecedents and the novelty boundary}
Integral freeness of the ordinary universal distribution is classical
\cite{Kubert}. Ouyang develops a broader universal norm distribution,
proves freeness, and gives an Anderson-type resolution
\cite[Propositions 4.1--4.2 and Theorem 5.1]{OuyangNorm}.
The ordinary sign-cohomology calculation is also classical;
Ouyang's spectral-sequence treatment records the result and its
antecedents \cite[Theorem 7.1]{OuyangSign}.

Accordingly, neither ordinary-distribution freeness nor the general idea
of an Anderson resolution is presented as a new discovery. The
contribution here is an explicit, self-contained treatment of the
manuscript's exact weighted presentation: a raw-row unit minor, an
original-grid-symbol compiler, the binary Koszul data for reflection,
and a closed minimum-valuation formula for arbitrary integral finite
jets. These statements are proved rather than inferred from an analogy
with the classical theory. Their precise extent of prior appearance in
the broader norm-distribution literature is not settled by the literature
check made for this report.

\subsection{Three different assertions}
A basis of a formal distribution quotient means independence modulo the
\emph{declared distribution rows}. It does not mean that the evaluated
complex numbers are linearly independent. An exact integer certificate
proves a relation from those rows; a small floating-point residual is
only a diagnostic. Finally, a torsion class in a reflected formal
quotient is not torsion of a complex period: evaluation in a
characteristic-zero vector space annihilates such a class.

These distinctions will be used throughout. In particular, none of the
rank or Smith calculations below proves irrationality of a Clausen value,
independence of Dirichlet $L$-values, or a minimal multiple-polylogarithm
depth.

\section{The weighted module and a fixed original-symbol basis}
\label{sec:presentation}
Let $q\ge2$ and write
\[
 q=\prod_{p\in P}p^{e_p},\qquad
 X_q=(q^{-1}\Z)/\Z,\qquad R=\Z[A_p:p\in P].
\]
All rings are commutative with identity. Put
\[
 M_q=\bigoplus_{x\in X_q}Re_x.
\]
For $p\in P$ and $x\in X_{q/p}$ define
\begin{equation}
 r_{p,x}=\sum_{\substack{y\in X_q\\py=x}}e_y-A_pe_x.
 \label{eq:raw}
\end{equation}
Let $\I_q$ be the $R$-span of every such row, and set
$\U_q=M_q/\I_q$. There are $q$ columns and
$m_q=\sum_{p\mid q}q/p$ raw rows. The symbol $e_0$ is retained.
For polylogarithmic evaluation it means $\Li_s(1)=\zeta(s)$;
for Hurwitz-zeta evaluation it means $\zeta(s,1)$, not $\zeta(s,0)$.

If $d\mid q$, put $A_d=\prod_p A_p^{v_p(d)}$.
The divisor row is defined in the same way, with $dy=x$ and weight $A_d$.
The identity
\begin{equation}
 r_{ab,x}=\sum_{ay=x}r_{b,y}+A_b r_{a,x}\qquad(ab\mid q)
 \label{eq:composite}
\end{equation}
proves by induction that prime rows generate all divisor rows. The
multiplicativity of composite weights is essential here.

\subsection{Canonical local selectors}
Write each point uniquely in its primary decomposition
\begin{equation}
 x=\sum_{p\in P}\frac{a_p(x)}{p^{e_p}},\qquad
 0\le a_p(x)<p^{e_p}.
 \label{eq:CRT}
\end{equation}
If $x=a/q$, the local numerator is
\[
 a_p(x)\equiv a\,(q/p^{e_p})^{-1}\pmod {p^{e_p}}.
\]
Define the local forbidden and retained sets
\begin{equation}
 T_p=\{1,\ldots,p^{e_p-1}\},\qquad
 B_p=\{0\}\cup\{p^{e_p-1}+1,\ldots,p^{e_p}-1\}.
 \label{eq:selectors}
\end{equation}
The proposed basis set is
\begin{equation}
 \B_q=\{x\in X_q:a_p(x)\in B_p\text{ for every }p\mid q\}.
 \label{eq:basis}
\end{equation}
It has $\varphi(q)$ elements and contains zero.
The set $T_p$ selects one \emph{nonzero} lift from each multiplication-by-$p$
fiber on the local grid. For a nonzero parent numerator
$b<p^{e_p-1}$ it selects $b$; for parent zero it selects
$p^{e_p-1}$.

Let $h_p(x)$ be the exponent of the exact order of the $p$-component of
$x$, with $h_p(x)=0$ for a zero component. Write
\[
 H(x)=\sum_p h_p(x),\qquad
 b(x)=\#\{p:a_p(x)\in T_p\}.
\]
For $x\notin\B_q$, let $p(x)$ be its least forbidden prime and choose
one original row
\begin{equation}
 g_x=r_{p(x),p(x)x}.
 \label{eq:selected}
\end{equation}
Order all grid points lexicographically by
$(H(x),b(x),a)$, where $x=a/q$ and $0\le a<q$.

\begin{theorem}[Original-symbol unimodularity]
\label{thm:unimodular}
The classes $[e_b]$, $b\in\B_q$, form an $R$-basis of $\U_q$.
The selected raw rows $g_x$, $x\notin\B_q$, form an $R$-basis of
$\I_q$. In the indicated order, their submatrix on the nonbasis columns
is unit lower triangular. In particular, its determinant is identically
one.

The complete $m_q\times q$ raw relation matrix is equivalent, by
invertible row and column operations over $R$, to an identity block of
size $q-\varphi(q)$ and zeros. All these statements survive arbitrary
base change $R\to T$. In particular,
\[
 \U_q\otimes_R T\simeq T^{\varphi(q)},\qquad
 (\U_q/Re_0)\otimes_R T\simeq T^{\varphi(q)-1}.
\]
\end{theorem}

\begin{example}[Small bases]
We identify $a/q$ with its numerator in the following display:
\[
\begin{array}{c|l}
 q& q\B_q\pmod q\\\hline
 3&0,2\\
 4&0,3\\
 6&0,4\\
 12&0,5,8,9\\
 15&0,1,4,6,7,9,10,12\\
 30&0,2,8,12,14,18,20,24.
\end{array}
\]
These are not primitive-root grids. Including lower-conductor symbols is
exactly what permits one basis to work for every weight without
exceptional denominators.
\end{example}

\section{Proof of the unit-minor theorem}
\label{sec:unitproof}

\subsection{Triangularity without division}
Fix a nonbasis point $x$, and put $p=p(x)$. In the row $g_x$, the point
$x$ has coefficient one. Its $p$-component is nonzero, and multiplication
by $p$ lowers its exact local order. Thus $px\ne x$ and the weight term
cannot change that pivot coefficient.

Every other point in $g_x$ precedes $x$. Indeed, the other solutions of
$py=px$ have the same components away from $p$. Among those of the same
$p$-order as $x$, the unique forbidden lift was $x$ itself; all the other
lifts are retained at $p$. Their height is unchanged and their forbidden
count decreases by one. A possible zero local lift has lower height.
Finally $px$ has height $H(x)-1$: multiplication by $p$ preserves the
orders at every other prime, even though it may change their forbidden
status. Height takes priority over forbidden count in the ordering.

Two different selected points cannot select the same row. For a fixed
prime and parent there is exactly one forbidden lift in its local fiber;
the other local components are uniquely fixed by the parent. Therefore
there are exactly $q-\varphi(q)$ distinct selected rows, and their
nonbasis submatrix is unit lower triangular.

Let $K_q$ be their span. Successive unit-pivot elimination gives a direct
sum
\begin{equation}
 M_q=K_q\oplus \bigoplus_{b\in\B_q}Re_b.
 \label{eq:split-selected}
\end{equation}
What remains is to prove that the unselected raw rows belong to $K_q$.
This step is necessary: a unit minor alone gives a lower bound on the
rank of the relation module, not its full description.

\subsection{The complete generic rank}
We give the characteristic-zero argument in full so that the integral
descent does not rely on unseen relation matrices. Extend scalars to a
field $F$ containing the fraction field of $R$ and all character values
of $(\Z/q\Z)^\times$. Split the grid according to exact denominator:
\[
 F[X_q]=\bigoplus_{m\mid q}F[(\Z/m\Z)^\times].
\]
Use $(\Z/1\Z)^\times$ as a one-element group. For a primitive character
$\chi$ of conductor $f\mid q$ and every $f\mid m\mid q$, put
\[
 v_{m,\chi}=\sum_{a\in(\Z/m\Z)^\times}\chi(a)e_{a/m}.
\]
The finite character transforms are invertible over $F$.

Fourier-transform the \emph{sources} of all $p$-rows, separately on their
exact-denominator strata. For $f\mid m$ and $pm\mid q$, the resulting
row is
\begin{equation}
 v_{pm,\chi}-\lambda_{p,m,\chi}v_{m,\chi},\qquad
 \lambda_{p,m,\chi}=\begin{cases}
 A_p,&p\mid m,\\ A_p-\chi(p),&p\nmid m.
 \end{cases}
 \label{eq:edges}
\end{equation}
When $p\mid m$, every preimage has denominator $pm$. When $p\nmid m$,
one preimage has denominator $m$; its character coefficient is
$\chi(p)\chi(b)$, since its numerator satisfies $pb\equiv a\pmod m$.
All remaining preimages have denominator $pm$. This proves
\eqref{eq:edges}. Because both source and target transforms are invertible,
these are all transformed rows, not a selected collection of identities.

Define
\begin{equation}
 E_{m,\chi}=
 \prod_{p\mid f}A_p^{v_p(m)-v_p(f)}
 \prod_{\substack{p\mid m\\p\nmid f}}
 A_p^{v_p(m)-1}(A_p-\chi(p)).
 \label{eq:conductor}
\end{equation}
The products of edge factors are independent of the order in which
primes are adjoined. Every $v_{m,\chi}$ reduces to
$E_{m,\chi}v_{f,\chi}$. Conversely, the map
$v_{m,\chi}\mapsto E_{m,\chi}$ annihilates every edge and sends
$v_{f,\chi}$ to one. The character block therefore has dimension exactly
one. There is one such block for each character modulo $q$, giving
\begin{equation}
 \dim_F(F[X_q]/F\I_q)=\varphi(q).
 \label{eq:generic-rank}
\end{equation}
This is the conductor argument of the earlier report, now used only
where the Fourier transform is legitimate \cite{repodistribution}.

\subsection{Integral descent and all-ring specialization}
The selected rows are independent and have the full generic rank
$q-\varphi(q)$. Hence $F K_q=F\I_q$. For any raw row $r$, its class in
$M_q/K_q$ becomes zero over $F$. But \eqref{eq:split-selected} identifies
this quotient with a free module over the integral domain $R$.
It injects into its scalar extension to $F$. Thus the class of $r$ was
already zero: $\I_q=K_q$.

This proves the basis statements. To reduce the entire matrix, put the
selected rows first, subtract their polynomial combinations from every
remaining row, and use the invertible triangular pivot block to clear
the nonbasis columns. A final column operation clears the retained
columns of the selected rows. Every operation is invertible over $R$;
no integer, weight, or character value is divided out. The result is
an identity block and zeros.

An identity between these polynomial matrices remains an identity under
any homomorphism $R\to T$. The decomposition, normal forms, and unit
minor consequently survive arbitrary specialization, including rings
with zero divisors and fields of characteristic dividing $q$ or
$\varphi(q)$. Since $0\in\B_q$, the anchored statement follows by
deleting one basis vector. This completes the proof of
Theorem~\ref{thm:unimodular}.

\section{Division-free certificates, size bounds, and unreflected jets}
\label{sec:compiler}
Write
\[
 [e_x]=\sum_{b\in\B_q}P_{x,b}(\boldsymbol A)[e_b],\qquad
 P_{x,b}\in\Z[\boldsymbol A].
\]
The proof gives the explicit recurrence
\begin{equation}
 \NF(x)=\begin{cases}
 e_x,&x\in\B_q,\\[2pt]
 A_p\NF(px)-\displaystyle\sum_{j=1}^{p-1}\NF(x+j/p),
   &x\notin\B_q,\ p=p(x).
 \end{cases}
 \label{eq:recurrence}
\end{equation}
Every argument on the second line precedes $x$. Thus this is a terminating
compiler over the polynomial ring itself, not a specialization-dependent
Gaussian elimination.

\begin{proposition}[Degree and coefficient bounds]
\label{prop:bounds}
For every $x,b$ and $p\mid q$,
\[
 \deg_{A_p}P_{x,b}\le h_p(x).
\]
Let $r=|P|$, let $P_{\max}$ be the largest prime divisor of $q$, and put
$W(x)=(r+1)H(x)+b(x)$. If the coefficient norm sums absolute values over
all basis coordinates and all monomials, then
\begin{equation}
 \|\NF(x)\|_1\le P_{\max}^{W(x)}.
 \label{eq:height}
\end{equation}
In particular each coordinate has at most
$\prod_p(h_p(x)+1)\le\prod_p(e_p+1)$ possible monomials.
\end{proposition}
\begin{proof}
A factor $A_p$ is introduced only on the branch $x\mapsto px$, which
lowers $h_p$ by one. All other local order exponents are unchanged.
Sibling lifts do not increase any local order. Induction in the
well-founded point order proves the multidegree bound.

Every recursion edge lowers $W$ by at least one. At equal height, the
forbidden count drops; when height drops, a possible increase of at most
$r$ in the forbidden count is outweighed by the factor $r+1$.
There are at most $p\le P_{\max}$ summands in a recursion step, including
the weighted parent. Multiplication by a variable does not increase the
coefficient norm. Induction proves \eqref{eq:height}.
\end{proof}

These estimates are deliberately uniform rather than sharp. Given the
factorization of $q$, straightforward sparse arithmetic uses at most
$O(qP_{\max}\varphi(q)\tau(q)r)$ integer coefficient operations under
this monomial bound, where $\tau(q)=\prod_p(e_p+1)$. This is a
polynomial-in-grid-size bound, not an assertion of optimal complexity
or of efficient factorization of a binary-encoded input integer.

\subsection{Independent replay}
A certificate records the basis, the selected original rows, and every
polynomial $P_{x,b}$. An independent verifier can check it without
re-running the discovery recurrence:
\begin{enumerate}[nosep]
\item Reconstruct each raw row directly from its prime and parent.
\item Verify that the selected nonbasis block is unit lower triangular.
\item Back-substitute $e_x-\sum_bP_{x,b}e_b$ through that block and require
an exactly zero polynomial remainder, for every $x$.
\item Independently require every raw row to vanish under the delivered
normal-form map.
\end{enumerate}
The delivered verifier follows this route and does not import the
normal-form generator. Altering a coefficient, the advertised basis, or
a selected parent is rejected by the corruption controls.

\begin{corollary}[Unreflected Smith factors and jets]
\label{cor:unreflected}
For every integer specialization of the prime weights, all nonzero Smith
factors of the complete raw matrix are one. Its rank is
$q-\varphi(q)$ over every field after the corresponding specialization.

For $B_L=\Z[u]/(u^L)$, $L\ge1$, and arbitrary weights $A_p(u)\in B_L$,
the quotient is free of rank $\varphi(q)$ over $B_L$. Expanded in the
coefficient basis $1,u,\ldots,u^{L-1}$, the raw integer matrix has
$L(q-\varphi(q))$ nonzero Smith factors, all equal to one.
\end{corollary}
\begin{proof}
Specialize the polynomial matrix equivalence from
Theorem~\ref{thm:unimodular}. For jets, multiplication by each entry is
represented in the coefficient basis. An invertible matrix over $B_L$
induces an invertible integer matrix on the underlying free abelian
group. Expanding the identity block gives the stated result.
\end{proof}

The same argument works for any finite free coefficient algebra, including
multivariate truncated polynomial algebras. It concerns ordinary
coefficients of formal jets; dividing derivatives by factorials is
unnecessary and may be illegitimate in positive characteristic.

\section{A weighted Anderson resolution with an elementary filtration proof}
\label{sec:resolution}
The reflection calculation requires relations among the raw relations,
not just their rank. We provide an explicit resolution. This is a
weighted version of the classical Anderson construction
\cite{OuyangNorm}; the exactness proof below is given for the present
point-grid convention.

For $S\subseteq P$, put $d_S=\prod_{p\in S}p$, and set
\begin{equation}
 \Lcal^{-k}=\bigoplus_{\substack{S\subseteq P\\|S|=k}}R[X_{q/d_S}].
 \label{eq:resolution-terms}
\end{equation}
Denote a generator by $[x,S]$. With the primes ordered increasingly,
define
\begin{equation}
 d[x,S]=\sum_{p\in S}(-1)^{\#\{\ell\in S:\ell<p\}}
 \left(\sum_{py=x}[y,S\setminus\{p\}]
       -A_p[x,S\setminus\{p\}]\right).
 \label{eq:resolution-d}
\end{equation}
Here the preimages lie in $X_{q/d_{S\setminus\{p\}}}$.
The augmentation $\Lcal^0\to\U_q$ is the defining quotient map.

\begin{theorem}[Exact weighted resolution]
\label{thm:resolution}
The augmented complex
\[
 0\longrightarrow\Lcal^{-|P|}\longrightarrow\cdots
 \longrightarrow\Lcal^{-1}\longrightarrow\Lcal^0
 \longrightarrow\U_q\longrightarrow0
\]
is exact and remains exact after arbitrary base change of $R$.
It is equivariant under $J[x,S]=[-x,S]$.
\end{theorem}
\begin{proof}
The unsigned operators in distinct prime directions commute: the two
iterated preimage sums give the same solutions of $p\ell y=x$, and
inclusion and the scalar weight terms commute with them. The alternating
signs in \eqref{eq:resolution-d} therefore give $d^2=0$. The augmentation
has cokernel zero and kernel $\I_q$ by definition. Equivariance is
immediate.

Filter the complex by
\[
 F(x,S)=\sum_p h_p(x)+|S|.
\]
Every differential term has filtration at most that of its source.
For a prime $p\in S$, the preimages of highest local order have
$p$-height $h_p(x)+1$. They preserve the filtration. A possible lower
local lift, and the scalar term $-A_p[x,S\setminus\{p\}]$, lower it.
In the associated graded complex, record
\[
 \alpha_p=h_p(x)+\boldsymbol1_{p\in S},\qquad0\le\alpha_p\le e_p.
\]
The leading differential preserves the whole tuple $\boldsymbol\alpha$.

There is a harmless prime-to-$p$ twist on the other primary components
of a preimage. Remove it in the summand indexed by $S$ by replacing
$x_p$ with
\begin{equation}
 \left(\prod_{\ell\in S\setminus\{p\}}\ell\right)^{-1}x_p.
 \label{eq:twist}
\end{equation}
The multiplier is a unit on the $p$-primary group. If the differential
removes $\ell\ne p$, its preimage operation multiplies $x_p$ by
$\ell^{-1}$, while the coordinate change in \eqref{eq:twist} gains the
factor $\ell$. These cancel. In the $p$-direction the multiplier is
unchanged. Thus the graded complex for a fixed
$\boldsymbol\alpha$ is a tensor product of independent local complexes.

For $\alpha_p=0$, the local complex is one copy of $R$ in degree zero.
For $\alpha_p\ge1$, it is a two-term complex from the free module on
points of exact order $p^{\alpha_p-1}$ to that on points of exact order
$p^{\alpha_p}$, sending a point to its highest-order preimage sum.
At $\alpha_p=1$, the source is the zero point and its image is the sum
of the $p-1$ nonzero $p$-torsion points. At larger $\alpha_p$, each
fiber has $p$ points. In either case the fibers are disjoint and
nonempty, so choosing one point per fiber gives a unit-pivot splitting.
The map is split injective with free cokernel.

Each local complex is consequently a direct sum of a contractible
identity-map complex and a free module in degree zero. Their tensor
product has no negative-degree homology. The filtration is finite, so
successively lifting a cycle's highest filtered part proves that the
original complex has no negative-degree homology either. This argument
uses only unit-pivot splittings and therefore works over every
coefficient ring. Its degree-zero homology is precisely $\U_q$.
\end{proof}

\begin{lemma}[Nonfixed support is preserved]
\label{lem:nonfixed}
In \eqref{eq:resolution-d}, a generator with $2x\ne0$ has no nonzero
coefficient on a generator $[y,T]$ with $2y=0$.
\end{lemma}
\begin{proof}
The weight term has the same point $x$. A fixed preimage $y$ would give
$2x=p(2y)=0$, a contradiction. Thus neither type of term can land on a
fixed point.
\end{proof}
This elementary support fact is the reason the reflection calculation
below collapses exactly; it cannot be replaced by an unsupported
assertion that a spectral sequence has no higher differentials.

\section{Reflection torsion as a binary Koszul calculation}
\label{sec:reflection}
Let $B$ be a commutative ring whose underlying abelian group is free,
let $a_p\in B$ be the specialized weights, and write
$U_B=\U_q\otimes_R B$. We allow $B$ to have infinite $\Z$-rank for
now. Reflection is the involution $Je_x=e_{-x}$. Define the two
\emph{coinvariant} quotients
\begin{equation}
 C_+(U_B)=U_B/(J-1)U_B,\qquad
 C_-(U_B)=U_B/(J+1)U_B.
 \label{eq:coinvariants}
\end{equation}
They are not the eigensublattices $\ker(J\mp1)$. Over a ring in which
two is invertible this distinction disappears after applying the usual
projectors; over $\Z$ it is essential.

\subsection{The two-periodic complex}
For a free abelian group $U$ with involution $J$, use the complex whose
differential from even degree is $J-1$ and from odd degree is $J+1$.
Its cohomology is denoted $\HH^n(J,U)$; in particular
\[
 \HH^0(J,U)=\frac{\ker(J-1)}{(J+1)U},\qquad
 \HH^1(J,U)=\frac{\ker(J+1)}{(J-1)U}.
\]
This is Tate cohomology for the two-element group, but these explicit
formulas suffice for its use here.

\begin{lemma}[Only elementary two-torsion]
\label{lem:torsion}
For a free abelian $U$,
\[
 \Tor_{\Z} C_+(U)\simeq\HH^1(J,U),\qquad
 \Tor_{\Z} C_-(U)\simeq\HH^0(J,U).
\]
Both torsion groups are killed by two.
\end{lemma}
\begin{proof}
If $n[v]=0$ in $U/(J-1)U$, apply $J+1$ to a representative equation.
Since $U$ has no integer torsion, $(J+1)v=0$. Then
$(J-1)v=-2v$, so $2[v]=0$. Conversely, every element of
$\ker(J+1)$ gives such a two-torsion class. This identifies the torsion
with $\HH^1$. Replacing $J$ by $-J$ gives the other assertion.
\end{proof}

Put $\overline B=B/2B$ and define
\begin{equation}
 r_* = \#\{p\mid q:p\text{ odd}\}+\boldsymbol1_{4\mid q}.
 \label{eq:effective-primes}
\end{equation}
The list of binary parameters is
\begin{equation}
 \boldsymbol f=
 \bigl(1-\overline a_p:\ p\mid q,\ p\text{ odd}\bigr),
 \quad\text{with }\overline a_2(1-\overline a_2)
       \text{ appended if }4\mid q.
 \label{eq:binary-parameters}
\end{equation}
Let $K_{\overline B}(\boldsymbol f)$ be the homological Koszul complex:
$K_k=\bigwedge^k\overline B^{r_*}$ and the differential contracts a
basis vector in direction $i$ by multiplication by $f_i$.

\begin{theorem}[Fixed-point Koszul reduction]
\label{thm:koszul}
For every integer $n$, there is an isomorphism of $\overline B$-modules
\begin{equation}
 \HH^n(J,U_B)\simeq
 \bigoplus_{\substack{0\le k\le r_*\\k\equiv n\ (2)}}
 H_k\bigl(K_{\overline B}(\boldsymbol f)\bigr).
 \label{eq:tate-koszul}
\end{equation}
The formula includes $q=2$, when the parameter list is empty.
\end{theorem}
\begin{proof}
Apply the two-periodic complex to each term of the exact resolution
in Theorem~\ref{thm:resolution}, after base change to $B$.
Use the usual sign on one differential so that the total differential
squares to zero. The horizontal range is finite. Computing horizontal
homology first therefore identifies the total cohomology with
$\HH^*(J,U_B)$; this is also obtained by a finite cycle-lifting argument
from the augmentation. No unbounded-filtration convergence assertion is
needed.

In each vertical column, separate the point basis into fixed points and
two-element orbits. Construct a subcomplex $D$ as follows. In even
vertical degree include all nonfixed point coordinates and twice each
fixed coordinate. In odd vertical degree include the entire column.
A nonfixed orbit has an acyclic two-periodic complex. A fixed coordinate
in $D$ has alternating modules $2B,B$ and differential $2:B\to2B$,
which is an isomorphism onto its target. Thus every vertical column of
$D$ is acyclic. The subcomplex is horizontally stable: by
Lemma~\ref{lem:nonfixed}, nonfixed source coordinates never produce
fixed coordinates, and twice a fixed coordinate still has even fixed
coefficients after taking the differential. Since the horizontal range
is finite, the total complex of $D$ is acyclic.

The quotient double complex is now particularly simple. Odd vertical
degrees are zero; even vertical degrees are the free $\overline B$-module
on the fixed points of the corresponding grid. All vertical
differentials vanish. Its total cohomology is the direct sum, with the
appropriate degree parities, of the homology of the horizontal
fixed-point complex.

A grid $X_m$ has just the fixed point $0$ if $m$ is odd, and has fixed
points $0,1/2$ if $m$ is even. In an odd-prime direction, each fixed
point has exactly one fixed preimage, so the induced map is multiplication
by $1-\overline a_p$. These give the odd-prime Koszul factors.

It remains to examine the direction $2$ when $q$ is even. If $v_2(q)=1$,
the fixed-point map is
\[
 \overline B\longrightarrow\overline B^2,\qquad
 1\longmapsto\binom{1-\overline a_2}{1}.
\]
It is split injective with one free cokernel coordinate; its contractible
summand may be removed. Thus it contributes no Koszul parameter.
If $v_2(q)\ge2$, the map on fixed points is
\begin{equation}
 \begin{pmatrix}
 1-\overline a_2&0\\
 1&\overline a_2
 \end{pmatrix}.
 \label{eq:two-matrix}
\end{equation}
The first column comes from the source zero, whose two fixed preimages
are $0$ and $1/2$. The source $1/2$ has only nonfixed preimages, so its
fixed contribution is the weight term. Swapping rows and using the
unit entry $1$ transforms \eqref{eq:two-matrix} into
$\operatorname{diag}(1,\overline a_2(1-\overline a_2))$ over
$\overline B$. Removing the identity-map summand leaves exactly the
additional Koszul parameter in \eqref{eq:binary-parameters}.

The other directions commute with these operations. The fixed-point
complex is therefore homotopy equivalent to
$K_{\overline B}(\boldsymbol f)$. In total degree $n$, a homological
degree $k$ term contributes precisely when $n+k$ is even. This proves
\eqref{eq:tate-koszul}.
\end{proof}

\subsection{Complete scalar Smith torsion}
\begin{theorem}[All integer prime weights]
\label{thm:scalar-torsion}
Let $q>2$ and specialize each $A_p$ to an arbitrary integer $a_p$.
Then, for either sign,
\begin{equation}
 C_\pm(U_{\Z})\simeq
 \Z^{\varphi(q)/2}\oplus(\Z/2\Z)^{h_q(\boldsymbol a)},
 \label{eq:scalar-Smith}
\end{equation}
where
\begin{equation}
 h_q(\boldsymbol a)=
 \begin{cases}
 0,&\text{some odd }p\mid q\text{ has }a_p\text{ even},\\
 2^{r_*-1},&\text{every odd-prime weight is odd}.
 \end{cases}
 \label{eq:scalar-h}
\end{equation}
For $q=2$, the two quotients are $C_+=\Z$ and $C_-=\Z/2\Z$.
\end{theorem}
\begin{proof}
Over $\F$, the parameter $\overline a_2(1-\overline a_2)$ is always zero.
If an odd-prime parameter is one, the Koszul complex has a contractible
factor and all its homology vanishes. Otherwise every parameter is zero,
so $H_k=\F^{\binom{r_*}{k}}$. For $q>2$ we have $r_*\ge1$, and the sum
of binomial coefficients of either parity is $2^{r_*-1}$.
Theorem~\ref{thm:koszul} and Lemma~\ref{lem:torsion} prove the torsion
formula.

To compute the free ranks without dividing an integral lattice, tensor
with $\Q$. The trace of $J$ on the quotient equals its alternating trace
on the finite resolution. A point-permutation module has trace equal
to its number of fixed points. For odd $q$ this alternating sum is
$(1-1)^{|P|}$. If $q=2m$ with $m$ odd, it is
$(2-1)(1-1)^{\omega(m)}$. If $4\mid q$, it has the factor $2-2$.
For every $q>2$ the trace is therefore zero. Since the total dimension is
$\varphi(q)$ and an involution is diagonalizable over $\Q$, its two
rational eigenspaces each have dimension $\varphi(q)/2$.
This is also the free rank of either coinvariant quotient. The
classification of finitely generated abelian groups proves
\eqref{eq:scalar-Smith}. At $q=2$, the unreflected quotient is generated
by $e_0$ and reflection is the identity.
\end{proof}

This result determines the complete nonzero Smith factors of the
\emph{stacked} distribution-and-reflection matrix: there are
$h_q$ factors equal to two and
$q-\varphi(q)/2-h_q$ factors equal to one. Its remaining
$\varphi(q)/2$ column directions are free. In characteristic two the
reflected quotient has dimension $\varphi(q)/2+h_q$; in every odd
characteristic it has dimension $\varphi(q)/2$.
The unreflected rank never changes.

\begin{example}[A torsion jump with constant unreflected rank]
For $q=3$, the single raw relation is
\[
 e_{1/3}+e_{2/3}=(a_3-1)e_0.
\]
At $a_3=1$, reflection is diagonal with signs $+1,-1$ in the basis
$e_0,e_{2/3}$, and both coinvariant quotients are
$\Z\oplus\Z/2\Z$. At $a_3=0$, both coinvariant quotients are
$\Z$. The unreflected lattice is $\Z^2$ in both cases.
For $q=60$, $r_*=3$; the torsion multiplicity is four when both $a_3$
and $a_5$ are odd, and zero otherwise, regardless of $a_2$.
\end{example}

The effective prime count is not optional. When $q=2m$ with $m$ odd,
the canonical bases for $m$ and $2m$ are the same rational points.
The natural inclusion identifies the two modules after adjoining $A_2$,
and respects reflection: all odd-prime fibers of a point in $X_m$ remain
inside $X_m$. Thus a lone factor of two must not double the ordinary
sign-cohomology multiplicity. For ordinary weights the theorem recovers
the classical result with this conductor convention
\cite{OuyangSign}.

\subsection{A universal cohomology module}
\begin{corollary}[Before weight specialization]
\label{cor:universal-tate}
Over $R=\Z[A_p:p\mid q]$, Tate cohomology in odd degree vanishes, while
in even degree it is
\[
 R\big/\bigl(2,\ 1-A_p\ (p\text{ odd}),\
 A_2(1-A_2)\ \text{if }4\mid q\bigr).
\]
A prime variable $A_2$ with $v_2(q)=1$ is not constrained in this quotient.
\end{corollary}
\begin{proof}
The binary parameters form a regular sequence in the polynomial ring
$R/2R$: the odd-prime linear factors eliminate distinct variables, and,
when present, $A_2(1-A_2)$ is a nonzero monic polynomial in the remaining
variable. Its Koszul homology is concentrated in degree zero and equals
the quotient by the parameter ideal. Theorem~\ref{thm:koszul} applies
because $R$ is free as an abelian group.
\end{proof}
The extra cohomology after scalar specialization is thus a real
base-change phenomenon. It cannot be read off by simply specializing
the already-computed universal cohomology module.

\section{An exact Smith formula for arbitrary finite integral jets}
\label{sec:jets}
Let
\[
 B_L=\Z[u]/(u^L),\quad L\ge1,\qquad
 a_p(u)\in B_L,
\]
and reduce coefficients modulo two. In
$\overline B_L=\F[u]/(u^L)$, form the parameters
\eqref{eq:binary-parameters}. For an element $f$ of this ring, define
\[
 \nu(f)=\min\{j:[u^j]f\ne0\},\qquad \nu(0)=L,
 \quad v=\min_{f\in\boldsymbol f}\nu(f).
\]
For $q>2$, the list is nonempty.

\begin{theorem}[Minimum-valuation jet formula]
\label{thm:jets}
For $q>2$, the two reflected jet quotients have the following Smith
decomposition as abelian groups:
\begin{equation}
 C_\pm(U_{B_L})\simeq
 \Z^{L\varphi(q)/2}\oplus(\Z/2\Z)^{v\,2^{r_*-1}}.
 \label{eq:jet-Smith}
\end{equation}
More precisely, their integer-torsion submodules are each isomorphic,
as $B_L$-modules, to
\begin{equation}
 \bigl(B_L/(2,u^v)\bigr)^{2^{r_*-1}}.
 \label{eq:jet-torsion-module}
\end{equation}
The direct-sum decomposition in \eqref{eq:jet-Smith} is a statement over
$\Z$, not a claim that the torsion-free part is free over $B_L$.
\end{theorem}
\begin{proof}
In the principal ideal ring $\overline B_L$, choose a parameter of least
valuation. It is a unit times $u^v$ unless every parameter is zero,
which is the same argument with $v=L$. A change of basis in the free
module defining the Koszul complex transforms the parameter list into
$(u^v,0,\ldots,0)$: first divide by its unit factor, then subtract its
multiples from the other entries. All these are invertible operations
over $\overline B_L$.

The complex is now the tensor product of the two-term complex given by
multiplication by $u^v$ and $r_*-1$ zero-differential exterior factors.
Both the kernel and the cokernel of multiplication by $u^v$ are
isomorphic to $\overline B_L/(u^v)$; the kernel is the ideal
$u^{L-v}\overline B_L$. Consequently
\[
 H_k\bigl(K_{\overline B_L}(\boldsymbol f)\bigr)
 \simeq\bigl(\overline B_L/(u^v)\bigr)^{\binom{r_*}{k}}.
\]
Summing either degree parity in Theorem~\ref{thm:koszul} proves
\eqref{eq:jet-torsion-module}. Each copy has $\F$-dimension $v$,
so Lemma~\ref{lem:torsion} gives the torsion multiplicity in
\eqref{eq:jet-Smith}.

The unreflected underlying lattice has rank $L\varphi(q)$.
The alternating fixed-point trace used in
Theorem~\ref{thm:scalar-torsion} is multiplied by $L$, and still
vanishes. Thus each rational eigenspace has dimension
$L\varphi(q)/2$. This proves the free rank and completes the Smith
calculation.
\end{proof}

\begin{example}[Five jet coefficients at level fifteen]
Take $L=5$, $a_3(u)=1+u^2$, and $a_5(u)=1+2u+u^3$.
The binary parameters are $u^2,u^3$, so $v=2$ and $r_*=2$. Therefore
\[
 C_\pm(U_{B_5})\simeq\Z^{20}\oplus(\Z/2\Z)^4.
\]
The constant jets $a_3=a_5=1$ would instead give ten two-torsion factors.
The constant terms of the weights agree; their higher binary
coefficients change the integral reflection structure.
\end{example}

\begin{example}[The two-adic jet parameter matters]
At $q=12$, $L=4$, let $a_2(u)=1+u$ and $a_3(u)=1+u^3$.
The two parameters are $u^3$ and $u+u^2$, not simply
$1-a_3$ and $1-a_2$. Thus $v=1$, $r_*=2$, and
\[
 C_\pm(U_{B_4})\simeq\Z^8\oplus(\Z/2\Z)^2.
\]
Although the parity of a \emph{scalar} two-prime weight does not affect
Theorem~\ref{thm:scalar-torsion}, a nonconstant two-prime jet can affect
Theorem~\ref{thm:jets}.
\end{example}

\subsection{Why the free abelian part need not be jet-free}
Take $q=3$, $L=2$, and $a_3=1+u$. In the unreflected basis
$e=e_0$, $f=e_{2/3}$, reflection satisfies
\[
 Je=e,\qquad Jf=ue-f.
\]
Hence
\begin{equation}
 C_+(U_{B_2})=B_2e\oplus B_2f\big/\langle ue-2f\rangle.
 \label{eq:nonflat}
\end{equation}
As an abelian group this has free generators $e,f$ and a torsion generator
$uf$ of order two. In its torsion-free quotient, multiplication by $u$
sends $e$ to $2f$ and $f$ to zero. This is not a free rank-one
$B_2$-module: the image of $u$ in a free rank-one module is a primitive
rank-one sublattice, whereas here the image is $2\Z f$.
Thus a claim of $B_L$-freeness for the reflected torsion-free part would
be false. By contrast, $C_-=B_2f\oplus(B_2/(2,u))e$ in this example.
Both signs have the same abelian Smith factors, but different module
extensions over the jet ring.

\subsection{Multivariate jets and arithmetic weights}
For a finite free $\Z$-algebra $B$ of rank $L$, the Koszul theorem remains
an exact answer. If $q>2$ and
$h_k=\dim_{\F} H_k(K_{B/2B}(\boldsymbol f))$, then
\[
 \Tor_\Z C_+\simeq(\Z/2\Z)^{\sum_{k\text{ odd}}h_k},\qquad
 \Tor_\Z C_-\simeq(\Z/2\Z)^{\sum_{k\text{ even}}h_k}.
\]
The two sums are equal: the Euler characteristic of the Koszul complex
is $L(1-1)^{r_*}=0$. A closed minimum-valuation description is special
to the one-variable principal ideal ring; it must not be transferred
unchanged to a multivariate Artinian ring.

For actual positive integral polylogarithm orders, $A_p=p^{1-n}$ usually
belongs to $\Z[1/q]$, rather than $\Z$. The corresponding coefficient
ring must be stated. If $q$ is even, inverting $q$ also inverts two and
removes the reflected two-torsion. For odd $q$, every $p^{1-n}$ is one
modulo two, and the ordinary maximal multiplicity persists over
$\Z[1/q]$. These are statements about the chosen formal coefficient
lattice, not about independence of the evaluated periods.

\section{Polylogarithm identities and pole-cancelled order derivatives}
\label{sec:polylog}
Put $\zeta_q=e^{2\pi i/q}$ and initially define
$\Li_s(z)=\sum_{n\ge1}z^n/n^s$ for the roots considered here and
$\Re s>1$. The standard continuation conventions are those of
\cite{DLMFpolylog,DLMFhurwitz}. Absolute convergence gives
\begin{equation}
 \sum_{py=x}\Li_s(e^{2\pi i y})
       =p^{1-s}\Li_s(e^{2\pi i x}).
 \label{eq:analytic-dist}
\end{equation}
Indeed the sum of $p$-th roots of unity filters the summation index to
multiples of $p$.

For a fixed rational angle, splitting the series into residue classes
also gives
\begin{equation}
 \Li_s(\zeta_q^a)=q^{-s}\sum_{r=1}^{q}\zeta_q^{ar}
                         \zeta(s,r/q).
 \label{eq:Hurwitz-Fourier}
\end{equation}
The Hurwitz-zeta terms have only a common simple pole at $s=1$, of
residue one. For $a\not\equiv0\pmod q$, their root-of-unity coefficients
sum to zero. Thus the left side is entire as a function of the order
$s$ at a nontrivial root; for $a=0$ it is $\zeta(s)$. This also provides
a direct continuation argument for all identities below.

\begin{theorem}[Fixed-basis identities for every complex order]
\label{thm:all-order}
For every grid index $a$, the compiler polynomials satisfy the
meromorphic identity
\begin{equation}
 \Li_s(\zeta_q^a)=
 \sum_{b\in q\B_q}P_{a,b}\bigl((p^{1-s})_{p\mid q}\bigr)
                         \Li_s(\zeta_q^b).
 \label{eq:polylog-normal}
\end{equation}
Here $q\B_q$ denotes the numerator representatives of the basis points.
The coefficient functions are finite exponential polynomials in $s$.
There are no exceptional weight denominators in this basis.
\end{theorem}
\begin{proof}
Every polynomial certificate is a linear combination of raw rows.
Evaluate it with $A_p=p^{1-s}$ and use \eqref{eq:analytic-dist} in the
absolutely convergent half-plane. Both sides continue meromorphically
by \eqref{eq:Hurwitz-Fourier}, so the identity theorem gives
\eqref{eq:polylog-normal} for every complex $s$.
\end{proof}
The theorem is an identity theorem for a fixed finite presentation, not
an assertion that these $\varphi(q)$ evaluated functions or any of their
special values are a minimal numerical basis.

\subsection{Two explicit all-order identities with short raw certificates}
At level twelve write $F_a(s)=\Li_s(\zeta_{12}^a)$. The normal form of
$F_1$ yields
\begin{equation}
 \boxed{\begin{aligned}
 &\Li_s(e^{i\pi/6})+\Li_s(e^{5i\pi/6})
       +(1+3^{1-s})\Li_s(-i)\\
 &\hspace{25mm}=(12^{1-s}-6^{1-s})\zeta(s).
 \end{aligned}}
 \label{eq:level12}
\end{equation}
In formal notation, with $\alpha=A_2$ and $\beta=A_3$, its certificate
is the sum of just three raw rows:
\begin{equation}
 r_{3,1/4}+\beta r_{2,1/2}+\alpha\beta r_{2,0}.
 \label{eq:level12-cert}
\end{equation}
Expanding gives
$e_{1/12}+e_{5/12}+(1+\beta)e_{3/4}
-(\alpha^2\beta-\alpha\beta)e_0$, as required.
Equivalently, ordinary triplication and duplication prove the same
formula. This familiar origin is not disguised as an independent new
functional equation; its role here is an all-order, pole-aware
integration test with an exact certificate.

At level thirty let $F_a(s)=\Li_s(\zeta_{30}^a)$, and put
$\alpha=2^{1-s}$, $\beta=3^{1-s}$, $\gamma=5^{1-s}$.
A less compact but useful generated identity is
\begin{equation}
\begin{aligned}
 F_1={}&(\beta\gamma-1)\zeta(s)+(\alpha-1)F_2
       -F_8-F_{12}-F_{14}\\
       &-(\beta+1)F_{18}-(\gamma+1)F_{20}-F_{24}.
\end{aligned}
\label{eq:level30}
\end{equation}
It too has a short raw proof. Before specialization, the residual of
\eqref{eq:level30} is exactly
\begin{equation}
 r_{2,1/15}-r_{3,3/5}+r_{5,1/3}+A_5r_{3,0}+r_{5,0}.
 \label{eq:level30-cert}
\end{equation}
For example, the first two rows cancel $e_{16/30}$, the next two
cancel $e_{26/30}$ and $e_{10/30}$, and the last cancels $e_{6/30}$.
The resulting endpoint coefficient is $1-A_3A_5$. Expanding these five
rows is a complete symbolic proof; no relation search is needed.

Both formulas specialize to identities for all positive integer
weights, to meromorphic values of the order, and, after cancellation,
to every order derivative at one. Appendix~\ref{app:q12} gives the entire
level-twelve polynomial normal-form table. The supplied JSON
certificates contain all point identities at seven levels.

\subsection{The endpoint coefficient vanishes at ordinary weights}
For a nonzero index $a$,
\begin{equation}
 P_{a,0}(1,\ldots,1)=0.
 \label{eq:zero-endpoint}
\end{equation}
There is a purely algebraic reason. At $A_p=1$, the zero-coordinate
coefficient in every raw row vanishes: for parent zero the two copies
of $e_0$ cancel, and for nonzero parent no zero preimage occurs.
Therefore every certified relation has zero coefficient on $e_0$.
Applying this to $e_a-\NF(a)$ proves \eqref{eq:zero-endpoint}.
The removable pole in \eqref{eq:polylog-normal} is consequently
certified before any numerical evaluation.

\begin{theorem}[All-order finite-part rule]
\label{thm:finitepart}
Let $a\ne0$, and define
\[
 C_{a,b}(t)=P_{a,b}((p^{-t})_p),\qquad
 L_b^{(j)}=\left.\frac{d^j}{ds^j}\Li_s(\zeta_q^b)\right|_{s=1}
 \quad(b\ne0).
\]
Use $\zeta(1+t)=t^{-1}+\sum_{k\ge0}(-1)^k\gamma_k t^k/k!$.
For every $j\ge0$,
\begin{align}
 L_a^{(j)}={}&
 \sum_{\substack{b\in q\B_q\\b\ne0}}
 \sum_{h=0}^{j}\binom jh C_{a,b}^{(h)}(0)L_b^{(j-h)}
 +\frac{C_{a,0}^{(j+1)}(0)}{j+1}\notag\\
 &+\sum_{h=1}^{j}\binom jh C_{a,0}^{(h)}(0)
                         (-1)^{j-h}\gamma_{j-h}.
 \label{eq:finitepart}
\end{align}
Coefficient derivatives are explicit. A monomial
$c\prod_p A_p^{m_p}$ contributes
$c(-\sum_pm_p\log p)^h$ to the $h$-th derivative, with the usual
order-zero convention.
\end{theorem}
\begin{proof}
Substitute $s=1+t$ in \eqref{eq:polylog-normal}.
Equation~\eqref{eq:zero-endpoint} gives $C_{a,0}(0)=0$.
Multiply its Taylor series by the displayed Laurent series of zeta.
The pole term contributes $C_{a,0}^{(j+1)}(0)/(j+1)$ after taking
$j$ derivatives; the regular zeta coefficients give the last sum.
The nonendpoint terms follow by Leibniz's rule. All series converge
in a sufficiently small punctured neighborhood, with a removable
singularity in the product, so coefficient extraction is legitimate.
\end{proof}
The derivative of order $j+1$ in the pole contribution is indispensable.
Replacing a removable product by a naive value-times-value expression
would lose this term.

For \eqref{eq:level12}, put
\[
 D_h=(-\log12)^h-(-\log6)^h,
 \quad F_a^{[j]}=\left.\partial_s^j\Li_s(\zeta_{12}^a)\right|_{s=1}.
\]
The complete derivative family becomes
\begin{align}
 F_1^{[j]}+F_5^{[j]}+F_9^{[j]}
 +\sum_{h=0}^{j}\binom jh(-\log3)^hF_9^{[j-h]}
 ={}&\frac{D_{j+1}}{j+1}\notag\\[-3pt]
 &\hspace{-45mm}+\sum_{h=1}^{j}\binom jhD_h(-1)^{j-h}\gamma_{j-h}.
 \label{eq:level12-jets}
\end{align}
In particular,
\begin{align}
 F_1^{[0]}+F_5^{[0]}+2F_9^{[0]}&=-\log2,\label{eq:jet0}\\
 F_1^{[1]}+F_5^{[1]}+2F_9^{[1]}-\log3\,F_9^{[0]}
 &=\frac{\log^2 12-\log^2 6}{2}-\gamma_0\log2.
 \label{eq:jet1}
\end{align}
For order two the right side is
\[
 \frac{-\log^3 12+\log^3 6}{3}
 +\gamma_0(\log^2 12-\log^2 6)+2\gamma_1\log2,
\]
and the left side is
$F_1^{[2]}+F_5^{[2]}+2F_9^{[2]}-2\log3\,F_9^{[1]}
+\log^2 3\,F_9^{[0]}$.

\subsection{Two exact cyclotomic product identities}
Using $\Li_1(z)=-\Log(1-z)$ with the branch obtained from the defining
disk and then continued to the indicated roots, \eqref{eq:jet0} gives
\begin{equation}
 (1-\zeta_{12})(1-\zeta_{12}^{5})(1-\zeta_{12}^{9})^2=2.
 \label{eq:product12}
\end{equation}
The pole-cancelled value of \eqref{eq:level30} gives
\begin{equation}
 \begin{aligned}
 &(1-\zeta_{30})(1-\zeta_{30}^{8})(1-\zeta_{30}^{12})
 (1-\zeta_{30}^{14})\\
 &\qquad\cdot(1-\zeta_{30}^{18})^2(1-\zeta_{30}^{20})^2
 (1-\zeta_{30}^{24})=15.
 \end{aligned}
 \label{eq:product30}
\end{equation}
Exponentiation of the exact logarithmic identities proves these products
without a branch ambiguity. Conversely, the product alone would not
fix the logarithmic identity's additive multiple of $2\pi i$.
As an independent algebraic check, the package gives integer polynomial
quotients verifying that each product polynomial minus its right side
is divisible by the corresponding cyclotomic polynomial.

\subsection{Stable numerical evaluation of the derivative identities}
Write the generalized Stieltjes expansion
\[
 \zeta(1+t,x)=t^{-1}+\sum_{k\ge0}\frac{(-1)^k\gamma_k(x)}{k!}t^k.
\]
Cancel the common pole in \eqref{eq:Hurwitz-Fourier} \emph{symbolically}.
For $a\ne0$ one obtains the finite identity
\begin{equation}
 \left.\partial_s^j\Li_s(\zeta_q^a)\right|_{s=1}
 =\frac{(-1)^j}{q}\sum_{k=0}^{j}\binom jk(\log q)^{j-k}
            \sum_{r=1}^{q}\zeta_q^{ar}\gamma_k(r/q).
 \label{eq:Stieltjes-evaluator}
\end{equation}
This follows by the same Laurent multiplication and is used in the
numerical script. It avoids evaluating a pole and subtracting it from
nearby values numerically. The numerical residuals remain diagnostics,
not proofs of the identities.

\section{Executable evidence and verification boundaries}
\label{sec:verification}
The proofs above are uniform in the level and jet length. The following
finite tests were actually run for the delivered package; the JSON
records contain individual cases, not just these totals.
\begin{center}
\small
\begin{tabular}{@{}p{.62\linewidth}r@{}}
\toprule
Exact check & Completed instances\\
\midrule
Polynomial normal forms, levels $2\le q\le120$ &119 levels\\
Raw polynomial rows annihilated on that range &3,317\\
Point normal forms on that range &7,259\\
Independent certificate files (levels $12,15,30,60,72,105,120$)&7\\
Point identities independently replayed in those files &414\\
Raw rows independently replayed in those files &366\\
Scalar parity patterns ($q\le500$, plus $840,1260,2310$)&2,348\\
Binary raw-row checks in the scalar sweep &488,954\\
Finite binary-jet cases &1,635\\
Direct raw reflected Smith computations &136\\
Direct raw unreflected Smith computations &68\\
Resolution cases over fields of sizes $2,3,5$ &498\\
Short displayed-identity raw-row certificates &2\\
Exact cyclotomic product polynomial divisions &2\\
Deliberately corrupted certificates rejected &3\\
\bottomrule
\end{tabular}
\end{center}

The scalar sweep includes \emph{every} binary prime-weight assignment at
each stated level. The jet sweep includes every binary coefficient
assignment of lengths one, two, and three at levels
$3,4,6,9,12,15,20,24,30,60$, and five further specified examples.
These finite ranges are not extrapolated as a substitute for the proofs.

The core generator and independent certificate verifier use only Python
integer arithmetic and the standard library. The direct Smith script
uses SymPy 1.14.0 and constructs the raw matrices without importing the
normal-form generator. The resolution checker likewise constructs the
complex directly from its rational grid labels. It checks $d^2=0$,
exactness in negative degrees over the tested fields, and the
nonfixed-support property of Lemma~\ref{lem:nonfixed}.

The numerical script uses mpmath and checks twenty cases at 70 and 110
decimal working digits: the level-$12$ and level-$30$ identities at
orders $2$, $5/2$, and $0.4+0.7i$, plus the level-$12$ derivative family
through order three at each precision. The ordinary-order left sides
use the finite Hurwitz Fourier sum, while the normal-form right sides
use polylogarithms. Derivatives use
\eqref{eq:Stieltjes-evaluator}. These computations use no interval
arithmetic; the saved residuals are not certified error bounds.

\subsection{Trust boundary of the artifacts}
The small independent verifier establishes the delivered polynomial
identities from reconstructed raw rows. Analytic validity of the raw
rows is proved in \eqref{eq:analytic-dist}; it is not an assumption
hidden in the verifier. The general homological results have written
proofs and finite implementation tests, not Lean or other
proof-assistant certificates. A checksum protects file integrity, not
mathematical validity. No repository files were changed during this
work.

\section{Audit corrections and proposed integration}
\label{sec:audit}
The existing characteristic-zero distribution theorem is not contradicted
by the results here. It is strengthened to an integral statement and
supplemented with reflection data. The following targeted changes address
wording still present in the pinned manuscript.

\subsection{Clausen weight parity versus conjugation parity}
In Chapter 2, the section on fundamental constants says that odd Clausen
values use odd characters. Under the chapter's declared convention,
$\Cl_{2m+1}(\theta)=\Re\Li_{2m+1}(e^{i\theta})$, an odd \emph{weight}
Clausen value is a cosine component and uses even characters. The
unambiguous replacement is:
\begin{quote}
Sine-series, or conjugation-odd, components use odd characters. With
$\Cl_{2m}=\Im\Li_{2m}$ and $\Cl_{2m+1}=\Re\Li_{2m+1}$, the odd-weight
Clausen components instead use even characters.
\end{quote}
This corrects a parity ambiguity without changing the subsequent
weight-two Clausen formulas.

\subsection{Failure of a reduction search is not nonreduction}
The Family 2 discussion in Chapter 4 states that the higher-weight
antisymmetric parts do not reduce against the $(i,1)$ basis alone.
The surrounding account concerns bounded integer-relation searches and
formal relation matrices. The numerical statement should instead say
that the reported search did not find those reductions. A formal
nonmembership statement must name the particular row presentation; it
must not be promoted to nonmembership of the evaluated periods in a
numerical algebra. This is consistent with the careful qualifications
already present elsewhere in the manuscript.

\subsection{Fourier transport is not unimodular lattice transport}
The manuscript's Fourier transform is exact over a cyclotomic field.
Its inverse contains group-order denominators, so equal dimensions and
transported vector-space relations do not by themselves identify Smith
forms of integer lattices. An integral use of that discussion should
specify the coefficient ring and any localization. Theorem~\ref{thm:unimodular}
provides a separate raw-row argument that avoids this obstruction for
the complete distribution presentation. This is a clarification of the
scope of Fourier transport, not a claimed error in the Fourier identity.

\subsection{Updating the research-status paragraph}
The scalar integral and modular questions, the reflected scalar Smith
forms, and the one-variable integral-jet Smith forms can now be marked
proved \emph{for this complete presentation}. Multivariate reflection
is reduced to an explicit Koszul homology problem, not to the
one-variable minimum-valuation formula. The full equivariant module
classification and numerical period-independence questions remain
separate. The integration fragment uses the label prefix
\code{intdist:} and is intended to accompany this report rather than
silently replace the earlier characteristic-zero proof.

\section{Further research questions}
\label{sec:future}
\paragraph{Equivariant integral module structure.}
Determine the full $R[J]$-module structure of the universal quotient and
of its reflected quotients, beyond Tate cohomology and abelian Smith
forms. The explicit jet example \eqref{eq:nonflat} shows that even the
torsion-free part has integral extension data that the Smith factors
miss. A useful first target is a classification of these extensions
at prime-power levels.

\paragraph{Multivariate jet torsion.}
For a specified Artinian algebra
$\F[u_1,\ldots,u_d]/\mathfrak a$, compute the Koszul homology of the
binary parameters in closed form or from a minimal resolution of their
ideal. Theorem~\ref{thm:koszul} gives the exact target. Any proposed
replacement for the minimum valuation must account for nonprincipal
ideals and their syzygies.

\paragraph{Optimizing the original-symbol basis.}
The forbidden lift in each local fiber, and the forbidden prime chosen
at a point, can be varied. Which choices minimize normal-form coefficient
height, support, or certificate length while retaining unit triangularity?
The present height bound is uniform but intentionally coarse.

\paragraph{Integral Fourier and character interfaces.}
Construct explicit comparison matrices between the original-symbol
basis, conductor-character coordinates, and primitive grids over the
smallest useful localized cyclotomic integer rings. The resulting
indices should separate coordinate singularities from actual torsion
and make integral transport between polylogarithm and Hurwitz-zeta
certificates auditable.

\paragraph{Higher-depth distributions.}
Extend the division-free certificate interface to depth-two distribution
relations together with specified shuffle, stuffle, and regularization
rows. Such a module is not the depth-one module treated here. The
surviving $S_6$ question requires additional analysis in that setting;
it is not resolved by reinterpreting a depth-one Smith computation.

\paragraph{Analytic jets with certified errors.}
Combine the exact finite-part formulas with interval-certified Hurwitz
or generalized Stieltjes evaluation, including explicit cancellation
budgets. This would turn the current diagnostics into certified
numerical enclosures while leaving equality proofs in the exact
polynomial layer.

\paragraph{Proof-assistant formalization.}
The shortest initial targets are the selector lemma, termination of
\eqref{eq:recurrence}, exact back-substitution, and the matrix
base-change theorem. A second stage can formalize the filtered
resolution and the acyclic-subcomplex proof of
Theorem~\ref{thm:koszul}. The analytic continuation layer can then be
attached independently.

\paragraph{Arithmetic evaluation kernels.}
Determine which relations among evaluated cyclotomic periods are not
visible in the declared distribution presentation. A formal basis or
nonzero formal class is not an independence theorem for numbers. Any
claim about the evaluation kernel should state its weight, coefficient
field, and available functional equations explicitly.

\clearpage
\appendix
\section{The complete level-twelve normal-form table}
\label{app:q12}
Let $E_a=[e_{a/12}]$, $\alpha=A_2$, $\beta=A_3$.
The retained basis is $E_0,E_5,E_8,E_9$. Each row gives the four
coefficients of $E_a$ in that basis.
\begin{center}
\renewcommand{\arraystretch}{1.28}
\begin{tabular}{@{}c r r r r@{}}
\toprule
$a$&$E_0$&$E_5$&$E_8$&$E_9$\\
\midrule
0&$1$&$0$&$0$&$0$\\
1&$\alpha\beta(\alpha-1)$&$-1$&$0$&$-\beta-1$\\
2&$\alpha(\beta-1)$&$0$&$-\alpha-1$&$0$\\
3&$\alpha(\alpha-1)$&$0$&$0$&$-1$\\
4&$\beta-1$&$0$&$-1$&$0$\\
5&$0$&$1$&$0$&$0$\\
6&$\alpha-1$&$0$&$0$&$0$\\
7&$\alpha(\beta-\alpha)$&$1$&$-\alpha(\alpha+1)$&$\beta+1$\\
8&$0$&$0$&$1$&$0$\\
9&$0$&$0$&$0$&$1$\\
10&$1-\beta$&$0$&$\alpha+1$&$0$\\
11&$\alpha(1-\beta)$&$-1$&$\alpha(\alpha+1)$&$0$\\
\bottomrule
\end{tabular}
\end{center}
The companion file \code{data/normal\_forms\_q12.json} contains the
unfactored integer monomial data. At arbitrary complex order, substitute
$\alpha=2^{1-s}$, $\beta=3^{1-s}$ and
$E_a=\Li_s(e^{2\pi ia/12})$. At formal integral jet weights, substitute
polynomials modulo $u^L$ instead. The same table supports both uses.

\section{Reproduction and artifact map}
\label{app:artifacts}
From the package root, the standard-library checks are
\begin{verbatim}
python code/run_checks.py
python code/check_resolution.py
python code/check_short_identities.py
python code/verify_certificates.py data/normal_forms_q12.json
\end{verbatim}
The optional independent matrix and numerical checks are
\begin{verbatim}
python -m pip install -r requirements-optional.txt
python code/check_smith.py
python code/check_products_and_controls.py
python code/numerical_checks.py
\end{verbatim}
The optional dependencies are SymPy and mpmath; the core verifier needs
neither. Build the article with \code{python code/build\_article.py} or
run pdfLaTeX twice in the \code{article} directory. Every data file names
its scope. In particular, \code{numerical\_checks.json} explicitly
states that it does not contain interval certificates.

The repository integration notes and focused corrections are in
\code{integration/}. The source snapshot and read file-blob hashes are
recorded in \code{PROVENANCE.json}. The checksum manifest is generated
only after compilation and final verification. No external font files
or third-party source PDFs are included.

\clearpage
\begin{thebibliography}{9}
\bibitem{repo}
V. Reshetnikov, ProveIt project,
\emph{Polylogarithms and their Arithmetic Bridges}, manuscript and
research-status chapters, repository snapshot \texttt{\snap}.
\url{https://github.com/VladimirReshetnikov/ProveIt/tree/28357e8ca63dd78327db91d9be239d75e4462879/Analysis/Polylogarithms/docs/manuscript}.

\bibitem{repodistribution}
ProveIt project, \emph{Distribution jets} research report,
\src{Analysis/Polylogarithms/docs/reports/distribution-jets/article/distribution_jets.tex},
repository snapshot as above; in particular the universal conductor
normal form and its characteristic-zero hypothesis.

\bibitem{Kubert}
D. S. Kubert, The universal ordinary distribution,
\emph{Bulletin de la Soci\'et\'e Math\'ematique de France}
\textbf{107} (1979), 179--202.
\href{https://doi.org/10.24033/bsmf.1891}{doi:10.24033/bsmf.1891}.

\bibitem{OuyangNorm}
Y. Ouyang, On the universal norm distribution,
\emph{Journal of the Ramanujan Mathematical Society}
\textbf{17}, no.~4 (2002), 287--311.
\href{https://arxiv.org/abs/math/0210297}{arXiv:math/0210297v2}.
The cited statements were checked in the arXiv version.

\bibitem{OuyangSign}
Y. Ouyang, Spectral sequences of universal distribution and Sinnott's
index formula, preprint (1999),
\href{https://arxiv.org/abs/math/9911031}{arXiv:math/9911031v1}.
See the fixed-point double-complex argument and Theorem~7.1.

\bibitem{DLMFhurwitz}
NIST Digital Library of Mathematical Functions,
Section 25.11, Hurwitz Zeta Function.
\url{https://dlmf.nist.gov/25.11}.

\bibitem{DLMFpolylog}
NIST Digital Library of Mathematical Functions,
Section 25.12, Polylogarithms.
\url{https://dlmf.nist.gov/25.12}.
\end{thebibliography}
\end{document}
